\scanpage{195}
% manuscript position: top
\mspage{111}
% manuscript position: 1
Je rappelle qu'on pouvait prendre
% manuscript position: 2
\[ I=\underline{\mathrm{Hom}}(J,C)_{/\varphi}\,, \]
% manuscript position: 3–4
où $\underline{\mathrm{Hom}}(J,C)\lhook\joinrel\longrightarrow\underline{\mathrm{Hom}}(J,\text{\deleted{$\unsure{\operatorname{Ind}_{\pi}|C}$}}\,M\deleted{\unsure{)\longrightarrow|)}}=\operatorname{Ind}'_{\pi_0}(C))$, quitte à \deleted{\unsure{ne}} pas exiger que $I$ soit ordonnée.

% manuscript position: 5–7
On a envie de prouver que $\operatorname{card}\mathrm{Fl}\,I\leq\pi$. L'ennui, c'est qu'on arrive seulement à la majoration
% manuscript position: 8
\[ \operatorname{card}\mathrm{Fl}\,I\leq\pi^{\pi_0}. \]
% manuscript position: 9–10
\deleted{(moyennant l'hypothèse $\pi\geq\pi_1^{\unsure{\pi_0|\pi}}$, \unsure{de f.|de fait} (c'est \unsure{(}} \deleted{\unsure{M}}
% manuscript position: 10–11
si $J$ est discrète, on n'a pas mieux, puisqu'alors
% manuscript position: 12
\[ I\simeq\prod_{j\in J} C_{/\varphi(j)}
\qquad\text{d'où }\mathrm{Fl}\,I=\prod_{j\in\unsure{I|J}}\mathrm{Fl}\,C_{/\varphi(j)}, \]
% manuscript position: 13–16
et si $\operatorname{Card}\mathrm{Fl}\,C_{/\varphi(j)}=\pi$ pour $\unsure{\forall|\mathrm{tt}}\unsure{,}\,j$, on trouve $\operatorname{card}\mathrm{Fl}\,I=\pi^{\pi_0}$ -- il n'y a pas moyen de faire mieux ! ]

% manuscript position: 17–19
Donc je n'arrive pas à prouver directement 4.5.4, ce qui m'aurait arrangé pour prouver la condition c) de loc. cit.
% manuscript position: 20–25
Mais, à vrai dire, \unsure{en|on} \deleted{\unsure{a vu déjà|a}} \unsure{rien que|déjà|une fois} par le travail de la démonstration du th. 4.3.6, partie $b\Rightarrow c$ -- on y a déjà prouvé que la filtration qu'on vient de construire là (par la condition a) \deleted{\unsure{que}} de 4.5.2), est une filtration cardinale, \deleted{La} la condition c) de SGA 4 I 9.12 est prouvée \unsure{à} la page 76. Mais je vais refaire
% manuscript position: foot
% archive sheet number: 194
